% =====================================================================
%  CHAPITRE 11 — GRAMMAIRE B2 : LA FINESSE
% =====================================================================
\chapter{Structures avancées : relatives, gérondif, comparaisons, emphase}

Vous savez former des phrases correctes. Ce chapitre vous apprend à former des phrases
\textbf{riches} --- celles qui font dire à un anglophone : « ton anglais est vraiment bon ».

\section{Unité 1 --- Les propositions relatives}

\begin{regle}[Les pronoms relatifs]
\en{who} (personnes, sujet) --- \en{whom} (personnes, complément, formel) ---
\en{whose} (possession) --- \en{which} (choses, animaux) --- \en{that} (personnes ou choses,
familier et très fréquent) --- \en{where} (lieu) --- \en{when} (temps) --- \en{why} (raison).
\end{regle}

\begin{tableau}[Relatives définissantes et non définissantes]
\begin{tabularx}{\linewidth}{@{}p{4.4cm}p{6.4cm}X@{}}
\entete{Type} & \entete{Règle} & \entete{Exemple} \tabularnewline
\hline
\textbf{Définissante} (essentielle) & Sans virgules. Elle identifie de qui ou de quoi on parle. \en{that} possible. Pronom complément omisssible. & \en{The man \textbf{who called} is my uncle.} / \en{The book \textbf{(that) I read} was great.} \tabularnewline
\textbf{Non définissante} (commentaire) & Entre virgules. Information supplémentaire. \en{that} \textbf{interdit}. Pronom jamais omis. & \en{My brother, \textbf{who lives in Dakar}, is a doctor.} \tabularnewline
\end{tabularx}
\end{tableau}

\begin{astuce}[Le pronom complément disparaît très souvent]
\en{The film (that) I saw yesterday} \ra \en{The film I saw yesterday}.
\en{The man (who) you met} \ra \en{The man you met}.
Cette omission est naturelle et fréquente --- mais impossible quand le pronom est sujet
(\en{*The man called is my uncle}).
\end{astuce}

\begin{tableau}[Les relatifs moins évidents]
\begin{tabularx}{\linewidth}{@{}p{3.4cm}p{4.4cm}X@{}}
\entete{Structure} & \entete{Emploi} & \entete{Exemple} \tabularnewline
\hline
\en{whose} & possession (personnes et choses) & \en{That's the girl \textbf{whose} father is a pilot.} \tabularnewline
\en{where} & lieu (= \en{in which}) & \en{This is the house \textbf{where} I was born.} \tabularnewline
\en{when} & temps (= \en{at which}) & \en{I remember the day \textbf{when} we met.} \tabularnewline
\en{why} & raison (après \en{reason}) & \en{That's the reason \textbf{why} I left.} \tabularnewline
\en{which} renvoyant à toute la phrase & commentaire sur ce qui précède & \en{He was late, \textbf{which} annoyed everyone.} \tabularnewline
\en{preposition + which/whom} & registre soutenu & \en{The city \textbf{in which} I live} = \en{The city \textbf{(that)} I live \textbf{in}} \tabularnewline
\end{tabularx}
\end{tableau}

\begin{piege}[Ne répétez pas le pronom]
\en{*The man who I met him yesterday} \ra \en{The man (who) I met yesterday}.
\en{*The book which I bought it} \ra \en{The book (which) I bought}.
Le relatif \emph{remplace} déjà le complément : en ajouter un second est la faute typique
des francophones qui traduisent mot à mot.
\end{piege}

\begin{exo}[1.1 --- Complétez avec \en{who}, \en{which}, \en{whose}, \en{where}, \en{when} ou rien]
\begin{enumerate}
  \item That's the woman \trou[1.5cm] helped me.
  \item The car \trou[1.5cm] I bought is blue.
  \item This is the school \trou[1.5cm] I studied.
  \item I'll never forget the day \trou[1.5cm] we met.
  \item He's the boy \trou[1.5cm] father is a minister.
  \item The film, \trou[1.5cm] was directed by a Senegalese, won a prize.
  \item The man \trou[1.5cm] you spoke to is my boss.
  \item She was late, \trou[1.5cm] made everyone angry.
\end{enumerate}
\end{exo}
\begin{corrige}[1.1]
1. \en{who / that} 2. \en{which / that / ---} (omis) 3. \en{where} 4. \en{when}
5. \en{whose} 6. \en{which} (non définissante : \en{that} interdit)
7. \en{who / whom / that / ---} 8. \en{which}
\end{corrige}

\section{Unité 2 --- Gérondif (\en{-ing}) ou infinitif (\en{to} + base) ?}

C'est l'un des points les plus subtils de l'anglais. Il existe trois situations.

\subsection{Cas 1 --- Verbes suivis toujours du gérondif}
\begin{tableau}[Verbes + \en{-ing} (liste essentielle)]
\footnotesize
\begin{tabularx}{\linewidth}{@{}X X X@{}}
\en{enjoy doing} & \en{avoid doing} & \en{finish doing} \\
\en{mind doing} & \en{keep doing} & \en{practise doing} \\
\en{suggest doing} & \en{consider doing} & \en{imagine doing} \\
\en{deny doing} & \en{admit doing} & \en{risk doing} \\
\en{miss doing} & \en{delay doing} & \en{can't help doing} \\
\en{give up doing} & \en{go on doing} & \en{feel like doing} \\
\en{it's no use doing} & \en{it's worth doing} & \en{look forward to doing} \\
\end{tabularx}
\end{tableau}
\ex{I enjoy learning English. She avoided answering. Do you mind waiting?}{J'aime apprendre l'anglais. Elle a évité de répondre. Ça t'ennuie d'attendre ?}
\ex{I'm looking forward to \textbf{seeing} you.}{J'attends de te voir avec impatience.}

\begin{piege}[\en{look forward to} + \en{-ing}]
Ici, \en{to} est une \textbf{préposition}, pas une marque d'infinitif : donc
\en{looking forward to meeting you} \ok, \en{*looking forward to meet you} \ko.
Même logique : \en{be used to doing}, \en{object to doing}, \en{dedicated to doing},
\en{when it comes to doing}.
\end{piege}

\subsection{Cas 2 --- Verbes suivis toujours de l'infinitif}
\begin{tableau}[Verbes + \en{to} + base (liste essentielle)]
\footnotesize
\begin{tabularx}{\linewidth}{@{}X X X@{}}
\en{want to do} & \en{decide to do} & \en{hope to do} \\
\en{plan to do} & \en{promise to do} & \en{refuse to do} \\
\en{agree to do} & \en{offer to do} & \en{manage to do} \\
\en{learn to do} & \en{would like to do} & \en{need to do} \\
\en{seem to do} & \en{appear to do} & \en{pretend to do} \\
\en{afford to do} & \en{choose to do} & \en{fail to do} \\
\en{expect to do} & \en{deserve to do} & \en{ask someone to do} \\
\end{tabularx}
\end{tableau}
\ex{I want to learn. He promised to come. We decided to leave.}{Je veux apprendre. Il a promis de venir. Nous avons décidé de partir.}

\subsection{Cas 3 --- Verbes qui changent de sens}
\begin{tableau}[Les huit verbes à double sens --- à connaître absolument]
\begin{tabularx}{\linewidth}{@{}p{2.4cm}p{4.2cm}p{4.2cm}X@{}}
\entete{Verbe} & \entete{+ \en{-ing}} & \entete{+ \en{to} + base} & \entete{Différence} \tabularnewline
\hline
\en{remember} & \en{I remember locking the door.} & \en{Remember to lock the door.} & souvenir du passé / tâche à faire \tabularnewline
\en{forget} & \en{I'll never forget meeting her.} & \en{Don't forget to call.} & souvenir / obligation \tabularnewline
\en{stop} & \en{He stopped smoking.} & \en{He stopped to smoke.} & arrêter une activité / s'arrêter pour faire autre chose \tabularnewline
\en{try} & \en{Try restarting the phone.} & \en{I tried to open it.} & expérimenter un moyen / faire un effort (souvent sans succès) \tabularnewline
\en{regret} & \en{I regret telling him.} & \en{I regret to inform you\ldots} & regretter un acte passé / annoncer poliment \tabularnewline
\en{go on} & \en{He went on talking.} & \en{He went on to explain.} & continuer la même chose / passer à autre chose \tabularnewline
\en{mean} & \en{Sorry, I didn't mean shouting.} & \en{I mean to finish tonight.} & --- / avoir l'intention \tabularnewline
\en{like} & \en{I like swimming.} (plaisir général) & \en{I like to swim in the morning.} (habitude choisie) & nuance très fine \tabularnewline
\end{tabularx}
\end{tableau}

\subsection{Après une préposition : toujours \en{-ing}}
\begin{regle}
Toute préposition (\en{about, after, before, for, of, without, instead of, because of,
interested in, good at, tired of, afraid of, sorry for, think about, dream of, succeed in,
believe in, insist on}) est suivie du gérondif.
\end{regle}
\ex{Thank you for coming. He left without saying goodbye. I'm interested in working with you.}{Merci d'être venu. Il est parti sans dire au revoir. Travailler avec vous m'intéresse.}

\subsection{L'infinitif de but et les autres constructions}
\begin{tableau}[Infinitif, \en{for}, \en{so that}]
\begin{tabularx}{\linewidth}{@{}p{4.2cm}X@{}}
\entete{Structure} & \entete{Exemple} \tabularnewline
\hline
\en{to} + base (but) & \en{I came \textbf{to help} you.} --- \textit{Je suis venu pour t'aider.} \tabularnewline
\en{in order to / so as to} (but, plus net) & \en{He studied hard \textbf{in order to pass}.} \tabularnewline
\en{so as not to} (but négatif) & \en{I whispered \textbf{so as not to wake} her.} \tabularnewline
\en{for} + nom (moyen) & \en{This knife is \textbf{for bread}.} \tabularnewline
\en{for} + personne + \en{to} + base & \en{It's hard \textbf{for me to understand}.} \tabularnewline
\en{so that} + proposition & \en{I called \textbf{so that} she \textbf{wouldn't} worry.} \tabularnewline
\en{too} + adjectif + \en{to} & \en{It's \textbf{too cold to swim}.} --- \textit{Trop froid pour nager.} \tabularnewline
\en{adjectif} + \en{enough} + \en{to} & \en{He's \textbf{old enough to drive}.} --- \textit{Assez grand pour conduire.} \tabularnewline
\end{tabularx}
\end{tableau}

\begin{piege}[\en{too} et \en{enough} : l'ordre change]
\en{too} \textbf{avant} l'adjectif : \en{too hot}. \en{enough} \textbf{après} l'adjectif :
\en{hot enough}. \en{*enough hot} est une faute immédiate et très visible.
\end{piege}

\begin{exo}[2.1 --- \en{-ing} ou \en{to} ?]
\begin{enumerate}
  \item I enjoy (listen) \trou[2cm] to music.
  \item She promised (help) \trou[2cm] me.
  \item Would you mind (open) \trou[2.5cm] the door?
  \item We decided (leave) \trou[2cm] early.
  \item He stopped (smoke) \trou[2cm] two years ago.
  \item He stopped (buy) \trou[2cm] some bread on his way home.
  \item Don't forget (send) \trou[2.5cm] the email.
  \item I'll never forget (visit) \trou[2cm] Paris.
  \item Thanks for (come) \trou[2cm] so quickly.
  \item It's worth (see) \trou[2cm].
  \item I'm looking forward to (hear) \trou[2cm] from you.
  \item She's very good at (explain) \trou[2cm] things.
  \item I tried (open) \trou[2cm] the window, but it was stuck.
  \item Try (press) \trou[2.5cm] the red button.
  \item I regret (say) \trou[2cm] those words.
\end{enumerate}
\end{exo}
\begin{corrige}[2.1]
1. \en{listening} 2. \en{to help} 3. \en{opening} 4. \en{to leave} 5. \en{smoking}
6. \en{to buy} 7. \en{to send} 8. \en{visiting} 9. \en{coming} 10. \en{seeing}
11. \en{hearing} 12. \en{explaining} 13. \en{to open} 14. \en{pressing}
15. \en{saying}
\end{corrige}

\section{Unité 3 --- Comparaison et superlatif}

\begin{tableau}[Formation régulière]
\begin{tabularx}{\linewidth}{@{}p{3.6cm}p{3.6cm}X@{}}
\entete{Type d'adjectif} & \entete{Comparatif} & \entete{Superlatif} \tabularnewline
\hline
1 syllabe & \en{-er} + \en{than} & \en{the} + \en{-est} \tabularnewline
& \en{taller than, faster than} & \en{the tallest, the fastest} \tabularnewline
1 syllabe, consonne-voyelle-consonne & on double & \en{bigger, the biggest} ; \en{hotter, the hottest} \tabularnewline
fin en \en{-e} & \en{-r} & \en{nicer, the nicest} \tabularnewline
consonne + \en{y} & \en{-ier} & \en{happier, the happiest} ; \en{easier, the easiest} \tabularnewline
2 syllabes et plus & \en{more} + adjectif + \en{than} & \en{the most} + adjectif \tabularnewline
& \en{more beautiful than} & \en{the most beautiful} \tabularnewline
Irréguliers & \en{good} \ra \en{better} ; \en{bad} \ra \en{worse} ; \en{far} \ra \en{further / farther} & \en{the best ; the worst ; the furthest} \tabularnewline
\end{tabularx}
\end{tableau}

\subsection{Les structures de comparaison utiles}
\begin{tableau}[Douze structures à maîtriser]
\begin{tabularx}{\linewidth}{@{}p{4.4cm}X@{}}
\entete{Structure} & \entete{Exemple et traduction} \tabularnewline
\hline
\en{as + adj + as} & \en{He is \textbf{as tall as} his father.} --- \textit{Il est aussi grand que son père.} \tabularnewline
\en{not as / not so + adj + as} & \en{It's \textbf{not as expensive as} I thought.} --- \textit{Ce n'est pas aussi cher que je pensais.} \tabularnewline
\en{much / far / a lot + comparatif} & \en{It's \textbf{much better}. She's \textbf{far more experienced}.} --- \textit{Bien meilleur. Bien plus expérimentée.} \tabularnewline
\en{a bit / slightly + comparatif} & \en{It's \textbf{a bit harder} than I expected.} --- \textit{Un peu plus difficile.} \tabularnewline
\en{the + comparatif\ldots, the + comparatif\ldots} & \en{\textbf{The more} you practise, \textbf{the better} you get.} --- \textit{Plus tu pratiques, mieux c'est.} \tabularnewline
\en{comparatif + and + comparatif} & \en{It's getting \textbf{colder and colder}.} --- \textit{Il fait de plus en plus froid.} \tabularnewline
\en{less + adjectif + than} & \en{This one is \textbf{less expensive than} that one.} \tabularnewline
\en{the least} & \en{the least interesting film} --- \textit{le film le moins intéressant} \tabularnewline
\en{the same as} & \en{Your phone is \textbf{the same as} mine.} --- \textit{Le même que le mien.} \tabularnewline
\en{different from} & \en{English is \textbf{different from} French.} (et non \en{*different than} en BrE) \tabularnewline
\en{as many / as much as} & \en{I don't eat \textbf{as much as} you.} / \en{There weren't \textbf{as many people as} last year.} \tabularnewline
\en{prefer A to B} & \en{I \textbf{prefer tea to coffee}.} / \en{I'd \textbf{rather stay than go}.} \tabularnewline
\end{tabularx}
\end{tableau}

\begin{piege}[\en{*more tall}, \en{*the most big}]
Un adjectif d'une syllabe prend \en{-er/-est}, jamais \en{more/most} :
\en{taller} \ok, \en{*more tall} \ko. Exception courante : on peut dire \en{more fun}
et \en{more clear} (acceptés), mais \en{funner} est évité.
\end{piege}

\subsection{Comparer des quantités et des verbes}
\ex{I earn more than you. She works harder than me. He drives more carefully than I do.}{Je gagne plus que toi. Elle travaille plus dur que moi. Il conduit plus prudemment que moi.}
\ex{There are fewer cars than yesterday but less traffic noise.}{Il y a moins de voitures qu'hier mais moins de bruit de circulation.}
\begin{memo}[\en{fewer} ou \en{less} ?]
\en{fewer} + dénombrable pluriel (\en{fewer people, fewer errors}) ;
\en{less} + indénombrable (\en{less water, less time, less money}).
\end{memo}

\begin{exo}[3.1 --- Comparez]
\begin{enumerate}
  \item Dakar is (big) \trou[2cm] than Thiès.
  \item This exercise is (difficult) \trou[2.5cm] than the last one.
  \item She is (good) \trou[2.5cm] at English than I am.
  \item It's (bad) \trou[2cm] day of my life.
  \item My car is not (expensive) \trou[3cm] as yours.
  \item The (far) \trou[2cm] you go, the (beautiful) \trou[2.5cm] it gets.
  \item He runs (fast) \trou[2.5cm] than his brother.
  \item I have (few / less) \trou[2cm] time than you.
  \item This is the (interesting) \trou[2.5cm] book I have ever read.
  \item The weather is getting (hot) \trou[2cm] and (hot) \trou[2cm].
\end{enumerate}
\end{exo}
\begin{corrige}[3.1]
1. \en{bigger} 2. \en{more difficult} 3. \en{better} 4. \en{the worst}
5. \en{as expensive} 6. \en{further} --- \en{more beautiful} 7. \en{faster}
8. \en{less} 9. \en{most interesting} 10. \en{hotter} --- \en{hotter}
\end{corrige}

\section{Unité 4 --- Les \en{question tags} (les petites questions en fin de phrase)}

\begin{regle}[Le mécanisme]
Phrase affirmative \ra tag \textbf{négatif} ; phrase négative \ra tag \textbf{affirmatif}.
On reprend \textbf{l'auxiliaire} de la phrase (ou \en{do/does/did} s'il n'y en a pas) et
le \textbf{pronom} correspondant au sujet.
\end{regle}

\begin{tableau}[Tags --- tableau complet]
\begin{tabularx}{\linewidth}{@{}p{5.4cm}X@{}}
\entete{Phrase} & \entete{Tag} \tabularnewline
\hline
\en{You're French,} & \en{aren't you?} \tabularnewline
\en{She isn't coming,} & \en{is she?} \tabularnewline
\en{They live here,} & \en{don't they?} \tabularnewline
\en{He works hard,} & \en{doesn't he?} \tabularnewline
\en{You went out,} & \en{didn't you?} \tabularnewline
\en{She can swim,} & \en{can't she?} \tabularnewline
\en{We should leave,} & \en{shouldn't we?} \tabularnewline
\en{You've finished,} & \en{haven't you?} \tabularnewline
\en{They'll come,} & \en{won't they?} \tabularnewline
\en{I'm late,} & \en{\textbf{aren't I}?} (irrégulier !) \tabularnewline
\en{Let's go,} & \en{\textbf{shall we}?} \tabularnewline
\en{Don't forget,} & \en{\textbf{will you}?} \tabularnewline
\en{Open the door,} & \en{\textbf{will you / can you}?} \tabularnewline
\en{Nobody called,} & \en{\textbf{did they}?} (sujet négatif \ra tag affirmatif, pronom \en{they}) \tabularnewline
\en{Everything is fine,} & \en{\textbf{isn't it}?} \tabularnewline
\end{tabularx}
\end{tableau}

\begin{astuce}[Deux intonations, deux sens]
Tag \textbf{descendant} \searrow\ = je suis presque sûr, je cherche simplement l'accord :
\en{It's hot today, isn't it?} \searrow\ (conversation normale).\\
Tag \textbf{montant} \nearrow\ = vraie question, je ne sais pas :
\en{You haven't seen my keys, have you?} \nearrow\\
C'est l'outil de \textbf{small talk} par excellence : \en{Nice weather, isn't it?}
\en{You're new here, aren't you?} \en{That was a great match, wasn't it?}
\end{astuce}

\section{Unité 5 --- \en{so}, \en{neither}, \en{either} : répondre en écho}

\begin{tableau}[Marquer l'accord et le désaccord]
\begin{tabularx}{\linewidth}{@{}p{5.0cm}X@{}}
\entete{Situation} & \entete{Réponse} \tabularnewline
\hline
\en{I like coffee.} (affirmatif) & \en{\textbf{So do I.}} --- \textit{Moi aussi.} \tabularnewline
\en{I'm tired.} (avec \en{be}) & \en{\textbf{So am I.}} \tabularnewline
\en{I can swim.} (avec modal) & \en{\textbf{So can I.}} \tabularnewline
\en{I've been there.} (avec \en{have}) & \en{\textbf{So have I.}} \tabularnewline
\en{I don't like tea.} (négatif) & \en{\textbf{Neither do I.}} / \en{Nor do I.} --- \textit{Moi non plus.} \tabularnewline
\en{I'm not hungry.} & \en{\textbf{Neither am I.}} \tabularnewline
\en{I can't come.} & \en{\textbf{Neither can I.}} \tabularnewline
Désaccord sur un affirmatif & \en{\textbf{I don't.}} / \en{Really? I do.} \tabularnewline
Désaccord sur un négatif & \en{\textbf{I do.}} / \en{Well, I am.} \tabularnewline
\en{either} en fin de négation & \en{I don't like it \textbf{either}.} --- \textit{Je ne l'aime pas non plus.} \tabularnewline
\end{tabularx}
\end{tableau}

\begin{piege}[\en{*Me too} après une négation]
\en{``I don't like fish.''} --- \en{*Me too.} \ko \ra \en{\textbf{Neither do I.}} ou
\en{\textbf{Me neither.}} (familier, très courant à l'oral).
\end{piege}

\section{Unité 6 --- Les quantifieurs avancés}

\begin{tableau}[Tous les quantifieurs, avec leur construction]
\footnotesize
\begin{longtable}{@{}p{3.0cm}p{2.6cm}p{4.4cm}p{3.6cm}@{}}
\entete{Quantifieur} & \entete{Avec} & \entete{Sens} & \entete{Exemple} \tabularnewline
\endfirsthead
\entete{Quantifieur} & \entete{Avec} & \entete{Sens} & \entete{Exemple} \tabularnewline
\endhead
\hline
\en{a few} & dénombrables & quelques (positif) & \en{I have a few friends here.} \tabularnewline
\en{few} & dénombrables & très peu (négatif) & \en{Few people know this.} \tabularnewline
\en{a little} & indénombrables & un peu (positif) & \en{I speak a little English.} \tabularnewline
\en{little} & indénombrables & très peu (négatif) & \en{There is little hope.} \tabularnewline
\en{several} & dénombrables & plusieurs & \en{several reasons} \tabularnewline
\en{many / much} & dénombr. / indénombr. & beaucoup (surtout négation, question) & \en{How many? How much?} \tabularnewline
\en{a lot of / lots of} & les deux & beaucoup (affirmatif naturel) & \en{a lot of money, a lot of cars} \tabularnewline
\en{plenty of} & les deux & largement assez & \en{There's plenty of time.} \tabularnewline
\en{enough} & les deux & assez & \en{enough food, enough chairs} \tabularnewline
\en{too much / too many} & indénombr. / dénombr. & trop & \en{too much noise, too many errors} \tabularnewline
\en{some / any} & les deux & du, des / (négation, question) & \en{some water, any idea?} \tabularnewline
\en{no / none} & --- & aucun & \en{I have no time.} / \en{None of us knew.} \tabularnewline
\en{all (of)} & les deux & tout, tous & \en{all the students, all of it} \tabularnewline
\en{both (of)} & deux & les deux & \en{both of my brothers} \tabularnewline
\en{either / neither} & deux & l'un ou l'autre / ni l'un ni l'autre & \en{either day is fine; neither is true} \tabularnewline
\en{each / every} & singulier & chaque / tous (sans exception) & \en{each student, every day} \tabularnewline
\en{most (of)} & les deux & la plupart & \en{most people, most of the time} \tabularnewline
\en{half of, a third of} & les deux & la moitié, un tiers & \en{half of the cake} \tabularnewline
\end{longtable}
\end{tableau}

\begin{memo}[\en{every} contre \en{each} contre \en{all}]
\en{Every day} = tous les jours, sans exception (vision d'ensemble).\\
\en{Each student} = chaque étudiant pris un par un (vision individuelle).\\
\en{All the students} = tous les étudiants (le groupe).\\
Tous trois se construisent avec un verbe au \textbf{singulier} pour \en{every} et \en{each} :
\en{Every child \textbf{has} a book.}
\end{memo}

\section{Unité 7 --- Les adverbes : place et formation}

\begin{tableau}[Formation et exceptions]
\begin{tabularx}{\linewidth}{@{}p{5.0cm}X@{}}
\entete{Règle} & \entete{Exemples} \tabularnewline
\hline
adjectif + \en{-ly} & \en{quick} \ra \en{quickly}, \en{careful} \ra \en{carefully}, \en{quiet} \ra \en{quietly} \tabularnewline
\en{-y} \ra \en{-ily} & \en{happy} \ra \en{happily}, \en{easy} \ra \en{easily}, \en{angry} \ra \en{angrily} \tabularnewline
\en{-le} \ra \en{-bly} & \en{terrible} \ra \en{terribly}, \en{simple} \ra \en{simply} \tabularnewline
\en{-ic} \ra \en{-ically} & \en{basic} \ra \en{basically}, \en{economic} \ra \en{economically} \tabularnewline
Irréguliers & \en{good} \ra \en{well} ; \en{fast} \ra \en{fast} ; \en{hard} \ra \en{hard} ; \en{late} \ra \en{late} \tabularnewline
Deux formes, deux sens & \en{hard} (avec effort) / \en{hardly} (presque pas) ; \en{late} (en retard) / \en{lately} (récemment) ; \en{near} (près) / \en{nearly} (presque) \tabularnewline
\end{tabularx}
\end{tableau}

\begin{piege}[\en{hardly} ne veut pas dire « durement »]
\en{He works \textbf{hard}.} = il travaille dur.\\
\en{He \textbf{hardly} works.} = il ne travaille presque pas.\\
Même piège : \en{lately} = récemment (pas « tardivement »), \en{nearly} = presque,
\en{eventually} = finalement (pas « éventuellement » !), \en{actually} = en fait (pas
« actuellement » !). Ces quatre adverbes sont des \textbf{faux amis} --- voir chapitre 13.
\end{piege}

\begin{tableau}[Place de l'adverbe]
\begin{tabularx}{\linewidth}{@{}p{3.4cm}p{4.4cm}X@{}}
\entete{Type} & \entete{Place} & \entete{Exemple} \tabularnewline
\hline
Fréquence & avant le verbe principal, après \en{be} & \en{She \textbf{often} travels. He is \textbf{always} late.} \tabularnewline
Manière & après le verbe ou l'objet & \en{He drove \textbf{carefully}. She speaks English \textbf{well}.} \tabularnewline
Lieu & après le verbe, avant le temps & \en{They waited \textbf{outside}.} \tabularnewline
Temps & en fin de phrase (ou en début pour insister) & \en{I saw him \textbf{yesterday}. \textbf{Yesterday} I saw him.} \tabularnewline
Degré & avant l'adjectif ou l'adverbe modifié & \en{\textbf{very} good, \textbf{really} quickly, \textbf{quite} interesting} \tabularnewline
Commentaire & en début de phrase & \en{\textbf{Unfortunately}, we lost. \textbf{Honestly}, I don't know.} \tabularnewline
\end{tabularx}
\end{tableau}

\begin{astuce}[L'ordre quand plusieurs adverbes coexistent]
\textbf{Manière \ra Lieu \ra Temps} (\en{MPT}) : \en{She sang \textbf{beautifully}
\textbf{at the concert} \textbf{last night}.} Mémorisez « MPT » et vous serez correct à
chaque fois.
\end{astuce}

\section{Unité 8 --- Mettre en relief : emphase, inversion, clivage}

\subsection{\en{do} emphatique}
\ex{I \textbf{do} like your dress! You \textbf{did} call me, didn't you? Do sit down.}{J'aime \emph{vraiment} ta robe ! Tu m'as \emph{bien} appelé, non ? Assieds-toi donc.}

\subsection{Les phrases clivées (\en{cleft sentences})}
\begin{tableau}[Trois structures pour insister]
\begin{tabularx}{\linewidth}{@{}p{4.4cm}X@{}}
\entete{Structure} & \entete{Exemple} \tabularnewline
\hline
\en{It is / was X that\ldots} & \en{\textbf{It was} Ali \textbf{that} broke the window.} --- \textit{C'est Ali qui a cassé la vitre.} \tabularnewline
\en{What I need is\ldots} & \en{\textbf{What I need} is a good rest.} --- \textit{Ce dont j'ai besoin, c'est de repos.} \tabularnewline
\en{The thing that\ldots / All I want is\ldots} & \en{\textbf{All I want} is some peace.} --- \textit{Tout ce que je veux, c'est la paix.} \tabularnewline
\en{What happened was that\ldots} & \en{\textbf{What happened was that} the power went off.} --- \textit{Ce qui s'est passé, c'est que le courant a sauté.} \tabularnewline
\end{tabularx}
\end{tableau}

\subsection{L'inversion après un adverbe négatif ou restrictif (registre soutenu)}
\begin{tableau}[Inversions à reconnaître et à oser]
\begin{tabularx}{\linewidth}{@{}p{4.6cm}X@{}}
\entete{Amorce} & \entete{Phrase} \tabularnewline
\hline
\en{Never} & \en{\textbf{Never have I} seen such a thing.} \tabularnewline
\en{Rarely / Seldom} & \en{\textbf{Rarely does} he arrive on time.} \tabularnewline
\en{Hardly\ldots when} & \en{\textbf{Hardly had} we left \textbf{when} it started to rain.} \tabularnewline
\en{No sooner\ldots than} & \en{\textbf{No sooner had} I sat down \textbf{than} the phone rang.} \tabularnewline
\en{Not only\ldots but also} & \en{\textbf{Not only is} she intelligent, \textbf{but} she \textbf{is also} kind.} \tabularnewline
\en{Under no circumstances} & \en{\textbf{Under no circumstances should} you open this door.} \tabularnewline
\en{Only after / Only when} & \en{\textbf{Only after} the meeting \textbf{did I} understand.} \tabularnewline
\en{Little} & \en{\textbf{Little did} he know what was coming.} \tabularnewline
\end{tabularx}
\end{tableau}

\begin{astuce}[Utilisez-en une seule, mais utilisez-la]
Dans un exposé ou un entretien, une inversion bien placée fait un effet remarquable :
\en{``Not only did we finish on time, but we also came in under budget.''}
(Nous avons non seulement fini à temps, mais en plus sous le budget.)
Une par discours suffit --- au-delà, cela devient affecté.
\end{astuce}

\section{Unité 9 --- Les articles de niveau avancé}

\begin{tableau}[Cas délicats, résolus]
\begin{tabularx}{\linewidth}{@{}p{5.4cm}X@{}}
\entete{Situation} & \entete{Article} \tabularnewline
\hline
Instruments de musique & \en{the} : \en{play \textbf{the} piano, the guitar} \tabularnewline
Sports et jeux & rien : \en{play football, chess, tennis} \tabularnewline
Repas & rien en général : \en{have breakfast} ; mais \en{a} si on le décrit : \en{We had \textbf{a} lovely dinner} \tabularnewline
Institutions au sens fonctionnel & rien : \en{go to school / to bed / to prison / to church / to hospital (BrE) / to work / home} \tabularnewline
Mêmes lieux, sens bâtiment & \en{the} : \en{I went to \textbf{the} school to meet the teacher.} \tabularnewline
Moyens de transport avec \en{by} & rien : \en{by car, by bus, by train, by plane} --- mais \en{on foot} \tabularnewline
Moyens avec un déterminant & \en{the / my} : \en{in \textbf{the} car, on \textbf{the} bus} \tabularnewline
Noms géographiques & rien : \en{Senegal, Africa, Lake Victoria, Mount Everest} --- \en{the} : \en{the Gambia, the USA, the Sahara, the Atlantic, the Alps, the Nile} \tabularnewline
Rues, places, aéroports, gares & rien : \en{Oxford Street, Times Square, Heathrow Airport, Victoria Station} \tabularnewline
\en{the} + adjectif = groupe & \en{the poor, the rich, the young, the unemployed} \tabularnewline
Noms abstraits généralisés & rien : \en{Life is hard. Love matters. Education counts.} \tabularnewline
\en{a} avec un sens de « par » & \en{twice \textbf{a} week, 50 kilometres \textbf{an} hour} \tabularnewline
\end{tabularx}
\end{tableau}

\section{Unité 10 --- Verbe + préposition : les combinaisons obligatoires}
En anglais, le choix de la préposition dépend du verbe, pas de la logique française.
Apprenez chaque bloc d'un seul tenant.

\begin{tableau}[Les 60 combinaisons verbe/préposition les plus utiles]
\footnotesize
\begin{longtable}{@{}p{4.6cm}p{4.6cm}p{4.6cm}@{}}
\entete{+ \en{about}} & \entete{+ \en{for}} & \entete{+ \en{to}} \tabularnewline
\endfirsthead
\entete{+ \en{about}} & \entete{+ \en{for}} & \entete{+ \en{to}} \tabularnewline
\endhead
\en{think about} & \en{wait for} & \en{listen to} \tabularnewline
\en{talk about} & \en{look for} & \en{belong to} \tabularnewline
\en{dream about / of} & \en{ask for} & \en{happen to} \tabularnewline
\en{worry about} & \en{pay for} & \en{speak to / with} \tabularnewline
\en{care about} & \en{apologise for} & \en{write to} \tabularnewline
\en{know about} & \en{hope for} & \en{reply to} \tabularnewline
\en{complain about} & \en{thank someone for} & \en{agree to / with} \tabularnewline
\en{argue about} & \en{blame someone for} & \en{object to} \tabularnewline
\entete{+ \en{at}} & \entete{+ \en{of}} & \en{look forward to} \tabularnewline
\en{look at} & \en{think of} & \en{used to} \tabularnewline
\en{laugh at} & \en{dream of} & \en{prefer A to B} \tabularnewline
\en{arrive at (lieu précis)} & \en{consist of} & \en{get used to} \tabularnewline
\en{be good at} & \en{be afraid of} & \en{introduce someone to} \tabularnewline
\en{be bad at} & \en{be proud of} & \entete{+ \en{with}} \tabularnewline
\en{be surprised at} & \en{be tired of} & \en{agree with (personne)} \tabularnewline
\en{shout at} & \en{be full of} & \en{deal with} \tabularnewline
\entete{+ \en{in}} & \en{be made of} & \en{argue with} \tabularnewline
\en{believe in} & \en{be aware of} & \en{live with} \tabularnewline
\en{arrive in (ville, pays)} & \en{die of} & \en{fill in / out} \tabularnewline
\en{succeed in} & \en{think of / about} & \en{compare with} \tabularnewline
\en{be interested in} & \en{remind someone of} & \en{help someone with} \tabularnewline
\en{take part in} & \entete{+ \en{from}} & \en{provide someone with} \tabularnewline
\en{spend money in / on} & \en{come from} & \en{be satisfied with} \tabularnewline
\en{result in} & \en{suffer from} & \en{be angry with (qqn)} \tabularnewline
\en{specialise in} & \en{differ from} & \en{be busy with} \tabularnewline
\en{be involved in} & \en{recover from} & \en{begin with} \tabularnewline
\en{depend on} & \en{protect someone from} & \en{stay with} \tabularnewline
\en{rely on} & \en{prevent someone from} & \en{share with} \tabularnewline
\en{insist on} & \en{graduate from} & \en{be pleased with} \tabularnewline
\en{congratulate someone on} & \en{hide from} & \en{work with} \tabularnewline
\en{work on} & \en{learn from} & \en{be careful with} \tabularnewline
\en{be based on} & \en{escape from} & \en{be strict with} \tabularnewline
\en{spend time on} & \en{be different from} & \en{be familiar with} \tabularnewline
\en{focus on} & \en{be absent from} & \en{be bored with} \tabularnewline
\en{be keen on} & \en{take something from} & \en{be honest with} \tabularnewline
\en{count on} & \en{be far from} & \en{be popular with} \tabularnewline
\end{longtable}
\end{tableau}

\begin{piege}[Trois prépositions que le français fait rater]
\en{*depend of} \ra \en{depend \textbf{on}} \qquad
\en{*listen something} \ra \en{listen \textbf{to}} something \qquad
\en{*discuss about} \ra \en{discuss something} (pas de préposition !)\\
Ajoutez : \en{*married with} \ra \en{married \textbf{to}} ; \en{*angry against} \ra
\en{angry \textbf{with}} (personne) / \en{angry \textbf{about}} (situation) ;
\en{*search something} \ra \en{search \textbf{for}} (ou \en{look for}).
\end{piege}

\begin{exo}[10.1 --- Prépositions]
Complétez : \en{about, at, for, from, in, of, on, to, with} (parfois rien) :
\begin{enumerate}
  \item It depends \trou[1.2cm] the weather.
  \item I'm waiting \trou[1.2cm] the bus.
  \item She's married \trou[1.2cm] a doctor.
  \item Let's discuss \trou[1.2cm] the problem tomorrow.
  \item I'm looking forward \trou[1.2cm] hearing from you.
  \item He's very good \trou[1.2cm] mathematics.
  \item I arrived \trou[1.2cm] Dakar \trou[1.2cm] Monday.
  \item Listen \trou[1.2cm] me, please.
  \item She's afraid \trou[1.2cm] spiders.
  \item I agree \trou[1.2cm] you completely.
  \item He was born \trou[1.2cm] a small village.
  \item Thanks \trou[1.2cm] your help.
\end{enumerate}
\end{exo}
\begin{corrige}[10.1]
1. \en{on} 2. \en{for} 3. \en{to} 4. rien 5. \en{to} 6. \en{at}
7. \en{in} --- \en{on} 8. \en{to} 9. \en{of} 10. \en{with} 11. \en{in} 12. \en{for}
\end{corrige}

\begin{exo}[Synthèse du chapitre 11 --- reformulez]
\begin{enumerate}
  \item \en{The woman is my aunt. She called you.} (avec \en{who})
  \item \en{I want to thank you. You helped me.} (avec \en{for})
  \item \en{It's too difficult. I can't understand it.} (avec \en{too\ldots to})
  \item \en{You practise a lot. You improve.} (avec \en{the more\ldots the more})
  \item \en{Ali broke the window.} (mise en relief)
  \item \en{I have never seen such a beautiful beach.} (inversion)
  \item \en{``I'm hungry.''} (rapporté : \en{she said})
  \item \en{She is tall. Her sister is tall too.} (avec \en{as\ldots as})
  \item \en{I don't like tea.} --- accord négatif
  \item \en{You are tired.} (avec un tag)
\end{enumerate}
\end{exo}
\begin{corrige}[Synthèse]
1. \en{The woman who called you is my aunt.}
2. \en{I want to thank you for helping me.}
3. \en{It's too difficult for me to understand.}
4. \en{The more you practise, the more you improve.}
5. \en{It was Ali who broke the window.}
6. \en{Never have I seen such a beautiful beach.}
7. \en{She said (that) she was hungry.}
8. \en{Her sister is as tall as she is.}
9. \en{Neither do I.} / \en{Me neither.}
10. \en{You're tired, aren't you?}
\end{corrige}

\begin{memo}[Les dix réflexes qui marquent le niveau B2]
\begin{enumerate}
  \item Une relative au lieu de deux phrases : \en{The man who called\ldots}
  \item \en{for + -ing} après un remerciement : \en{Thanks for coming.}
  \item \en{look forward to + -ing} dans chaque conclusion d'e-mail.
  \item \en{used to} pour les souvenirs : \en{I used to live\ldots}
  \item \en{should have} pour les regrets : \en{I should have told you.}
  \item Un tag en fin de phrase pour créer du lien : \en{It's nice, isn't it?}
  \item \en{So do I / Neither do I} au lieu de \en{me too / me neither} systématiques.
  \item \en{as\ldots as} et \en{much better than} pour nuancer.
  \item \en{It was X who\ldots} pour mettre en relief.
  \item Le passif quand l'auteur importe peu : \en{The meeting was cancelled.}
\end{enumerate}
\end{memo}
